\chapter{Manual de Usuario}\label{manual-de-usuario}

\textbf{OrbitEngine} --- Plataforma SaaS para la Gestión de Procesos Internos en Pymes\\
Versión: 1.0 \textbar{} Abril 2026\\
Universidad Sergio Arboleda --- Semillero de Software como Innovación

\begin{center}\rule{0.5\linewidth}{0.5pt}\end{center}

\section{Introducción}\label{introducciuxf3n}

Este manual describe el uso de OrbitEngine para los usuarios finales de la plataforma. OrbitEngine es una aplicación web multi-tenant que permite a pequeñas y medianas empresas gestionar en un solo lugar sus operaciones de \textbf{inventario}, \textbf{ventas} y \textbf{clientes}, con un panel de indicadores clave de desempeño (KPIs) accesible en tiempo real.

El sistema opera completamente desde el navegador web, sin necesidad de instalar software adicional. Cada empresa (organización) tiene su propio espacio de trabajo completamente aislado de las demás.

\section{Acceso a la Plataforma}\label{acceso-a-la-plataforma}

\subsection{Página de inicio}\label{puxe1gina-de-inicio}

Al ingresar a la URL de OrbitEngine, el usuario accede a la \textbf{página de presentación} (landing page) que muestra las características de la plataforma. Desde esta página se puede navegar a las opciones de inicio de sesión o registro.

\subsection{Registro de organización}\label{registro-de-organizaciuxf3n}

Para crear una nueva empresa en la plataforma:

\begin{enumerate}
\def\labelenumi{\arabic{enumi}.}
\tightlist
\item
  Hacer clic en \textbf{``Comenzar gratis''} o en el enlace de registro en la barra de navegación.
\item
  En la pantalla \texttt{/signup-org}, completar el formulario con:

  \begin{itemize}
  \tightlist
  \item
    \textbf{Nombre de la organización} --- nombre visible de la empresa.
  \item
    \textbf{Slug} --- identificador único en minúsculas (ej. \texttt{mi-empresa}). Solo letras, números y guiones; entre 3 y 50 caracteres.
  \item
    \textbf{Descripción} \emph{(opcional)} --- breve descripción del negocio.
  \item
    \textbf{Nombre}, \textbf{apellido} y \textbf{correo electrónico} del administrador.
  \item
    \textbf{Contraseña} de al menos 8 caracteres.
  \end{itemize}
\item
  Hacer clic en \textbf{``Crear organización''}.
\item
  El sistema crea la organización y la cuenta administrador automáticamente y redirige al panel de control.
\end{enumerate}

\subsection{Invitación de usuarios}\label{invitaciuxf3n-de-usuarios}

Los usuarios adicionales \textbf{no se registran por su cuenta}: son creados por el administrador desde el módulo \textbf{Admin}. Una vez creada la cuenta, el usuario recibe las credenciales y puede ingresar con ellas.

\subsection{Inicio de sesión}\label{inicio-de-sesiuxf3n}

\begin{enumerate}
\def\labelenumi{\arabic{enumi}.}
\tightlist
\item
  Ir a la pantalla \texttt{/login}.
\item
  Ingresar el \textbf{correo electrónico} y la \textbf{contraseña}.
\item
  Hacer clic en \textbf{``Ingresar''}.
\item
  Si las credenciales son correctas, el sistema redirige al \textbf{panel de control} (\texttt{/dashboard}).
\end{enumerate}

\subsection{Recuperación de contraseña}\label{recuperaciuxf3n-de-contraseuxf1a}

Si olvidó su contraseña:

\begin{enumerate}
\def\labelenumi{\arabic{enumi}.}
\tightlist
\item
  En la pantalla de inicio de sesión, hacer clic en \textbf{``¿Olvidaste tu contraseña?''}.
\item
  Ingresar el correo electrónico asociado a la cuenta.
\item
  Revisar la bandeja de entrada y hacer clic en el enlace recibido.
\item
  En la pantalla \texttt{/reset-password}, ingresar y confirmar la nueva contraseña.
\item
  Hacer clic en \textbf{``Restablecer contraseña''}.
\end{enumerate}

\subsection{Cierre de sesión}\label{cierre-de-sesiuxf3n}

Para cerrar sesión, hacer clic en el \textbf{ícono de usuario} en la barra lateral izquierda y seleccionar \textbf{``Cerrar sesión''}. La sesión expira automáticamente después de 8 días de inactividad.

\section{Navegación General}\label{navegaciuxf3n-general}

Una vez dentro del sistema, la interfaz se divide en:

\begin{itemize}
\tightlist
\item
  \textbf{Barra lateral izquierda} --- menú principal de navegación con acceso a todos los módulos.
\item
  \textbf{Área de contenido principal} --- donde se muestra y gestiona la información de cada módulo.
\item
  \textbf{Panel de usuario} (parte inferior de la barra lateral) --- acceso a configuración personal y cierre de sesión.
\end{itemize}

Los módulos disponibles son:

{\def\LTcaptype{none} % do not increment counter
\begin{longtable}[]{@{}lll@{}}
\toprule\noalign{}
Módulo & Ruta & Acceso \\
\midrule\noalign{}
\endhead
\bottomrule\noalign{}
\endlastfoot
Dashboard & \texttt{/dashboard} & Todos los usuarios \\
Inventario & \texttt{/dashboard/inventory} & Todos los usuarios \\
Ventas & \texttt{/dashboard/sales} & Todos los usuarios \\
Clientes & \texttt{/dashboard/customers} & Todos los usuarios \\
Administración & \texttt{/dashboard/admin} & Solo Administradores \\
Configuración & \texttt{/dashboard/settings} & Todos los usuarios \\
\end{longtable}
}

\section{Dashboard --- Panel de Indicadores}\label{dashboard-panel-de-indicadores}

El dashboard muestra un resumen del estado del negocio en tiempo real para la organización activa.

\textbf{Indicadores disponibles:}

\begin{itemize}
\tightlist
\item
  \textbf{Ventas del día} --- total en dinero de las ventas registradas hoy.
\item
  \textbf{Ventas del mes} --- total acumulado del mes en curso.
\item
  \textbf{Ticket promedio} --- valor promedio por venta en el período.
\item
  \textbf{Productos con stock bajo} --- conteo de productos que están por debajo del mínimo configurado.
\item
  \textbf{Top productos} --- listado de los productos más vendidos.
\item
  \textbf{Ventas por día} --- gráfica de barras con la evolución diaria de ventas.
\end{itemize}

El dashboard se actualiza automáticamente con cada consulta. No requiere acción adicional del usuario.

\section{Módulo de Inventario}\label{muxf3dulo-de-inventario}

El módulo de inventario permite gestionar los productos y categorías de la empresa, así como llevar el control detallado de los movimientos de stock.

\subsection{Categorías}\label{categoruxedas}

Las categorías permiten organizar el inventario en grupos lógicos. Admiten \textbf{jerarquía} (una categoría puede tener subcategorías).

\textbf{Crear una categoría:}

\begin{enumerate}
\def\labelenumi{\arabic{enumi}.}
\tightlist
\item
  En el módulo Inventario, ir a la pestaña \textbf{Categorías}.
\item
  Hacer clic en \textbf{``Agregar categoría''}.
\item
  Completar el \textbf{nombre} \emph{(requerido)} y opcionalmente una \textbf{descripción} y una \textbf{categoría padre}.
\item
  Guardar.
\end{enumerate}

\textbf{Editar o desactivar una categoría:}

\begin{itemize}
\tightlist
\item
  Hacer clic en el menú de acciones (⋮) de la fila y seleccionar \textbf{``Editar''} o \textbf{``Desactivar''}.
\end{itemize}

\subsection{Productos}\label{productos}

\textbf{Crear un producto:}

\begin{enumerate}
\def\labelenumi{\arabic{enumi}.}
\tightlist
\item
  En el módulo Inventario, ir a la pestaña \textbf{Productos}.
\item
  Hacer clic en \textbf{``Agregar producto''}.
\item
  Completar el formulario:

  \begin{itemize}
  \tightlist
  \item
    \textbf{Nombre} \emph{(requerido)}
  \item
    \textbf{SKU} --- código único del producto \emph{(requerido)}
  \item
    \textbf{Categoría} \emph{(opcional)}
  \item
    \textbf{Precio de costo} y \textbf{precio de venta}
  \item
    \textbf{Stock actual}, \textbf{stock mínimo} y \textbf{stock máximo} \emph{(opcional)}
  \item
    \textbf{Unidad de medida} (ej. \texttt{unit}, \texttt{kg}, \texttt{lt})
  \item
    \textbf{Código de barras} \emph{(opcional)}
  \item
    \textbf{Descripción} e \textbf{imagen} \emph{(opcionales)}
  \end{itemize}
\item
  Guardar.
\end{enumerate}

\textbf{Editar un producto:}

\begin{itemize}
\tightlist
\item
  Hacer clic en el menú de acciones (⋮) → \textbf{``Editar''}. Todos los campos son modificables excepto el SKU si ya tiene movimientos.
\end{itemize}

\textbf{Desactivar un producto:}

\begin{itemize}
\tightlist
\item
  Hacer clic en el menú de acciones (⋮) → \textbf{``Desactivar''}. El producto deja de aparecer en los formularios de venta pero su historial se conserva.
\end{itemize}

\textbf{Alerta de stock bajo:}

\begin{itemize}
\tightlist
\item
  Los productos con \texttt{stock\_quantity\ ≤\ stock\_min} aparecen resaltados en la tabla y se contabilizan en el indicador del dashboard.
\end{itemize}

\subsection{Ajuste de stock manual}\label{ajuste-de-stock-manual}

Para corregir el inventario por diferencias de conteo, mermas u otros motivos:

\begin{enumerate}
\def\labelenumi{\arabic{enumi}.}
\tightlist
\item
  En la tabla de productos, seleccionar el producto y hacer clic en \textbf{``Ajustar stock''}.
\item
  Ingresar la \textbf{cantidad} (positiva para agregar, negativa para restar).
\item
  Ingresar un \textbf{motivo} del ajuste \emph{(requerido)}.
\item
  Confirmar.
\end{enumerate}

El ajuste queda registrado en el historial de movimientos con el tipo \texttt{adjustment}.

\subsection{Historial de movimientos}\label{historial-de-movimientos}

El historial registra todos los cambios de stock de cada producto:

{\def\LTcaptype{none} % do not increment counter
\begin{longtable}[]{@{}ll@{}}
\toprule\noalign{}
Tipo de movimiento & Descripción \\
\midrule\noalign{}
\endhead
\bottomrule\noalign{}
\endlastfoot
\texttt{sale} & Salida generada por una venta \\
\texttt{adjustment} & Entrada o salida manual \\
\texttt{return} & Entrada por devolución o cancelación de venta \\
\texttt{purchase} & Entrada por compra \\
\end{longtable}
}

Para ver el historial de un producto:

\begin{enumerate}
\def\labelenumi{\arabic{enumi}.}
\tightlist
\item
  En la tabla de productos, hacer clic en el menú de acciones (⋮) → \textbf{``Ver historial de movimientos''}.
\item
  Se muestra la tabla con fecha, tipo, cantidad, stock anterior y stock resultante, y el usuario que realizó el movimiento.
\end{enumerate}

\section{Módulo de Ventas}\label{muxf3dulo-de-ventas}

El módulo de ventas permite registrar transacciones comerciales, visualizar el historial y cancelar ventas cuando sea necesario.

\subsection{Registrar una venta}\label{registrar-una-venta}

\begin{enumerate}
\def\labelenumi{\arabic{enumi}.}
\tightlist
\item
  En el módulo Ventas, hacer clic en \textbf{``Nueva venta''}.
\item
  Completar el formulario:

  \begin{itemize}
  \tightlist
  \item
    \textbf{Cliente} \emph{(opcional)} --- buscar por nombre o documento.
  \item
    \textbf{Método de pago}: Efectivo, Tarjeta, Transferencia u Otro.
  \item
    \textbf{Descuento} y \textbf{IVA/impuesto} \emph{(opcional, en valores monetarios)}.
  \item
    \textbf{Notas} \emph{(opcional)}.
  \end{itemize}
\item
  En la sección de \textbf{ítems}, agregar los productos de la venta:

  \begin{itemize}
  \tightlist
  \item
    Buscar el producto por nombre o SKU.
  \item
    Ingresar la \textbf{cantidad}.
  \item
    El sistema calcula el precio unitario y el subtotal del ítem automáticamente.
  \end{itemize}
\item
  El sistema calcula en tiempo real: subtotal, descuento, impuesto y \textbf{total}.
\item
  Hacer clic en \textbf{``Registrar venta''}.
\end{enumerate}

Al guardar, el sistema:

\begin{itemize}
\tightlist
\item
  Genera un \textbf{número de factura} único.
\item
  Descuenta las cantidades del stock de cada producto registrado.
\item
  Registra los movimientos de inventario correspondientes (\texttt{sale}).
\item
  Actualiza el total de compras del cliente, si se asoció uno.
\end{itemize}

\subsection{Ver detalle de una venta}\label{ver-detalle-de-una-venta}

En la tabla de ventas, hacer clic en el menú de acciones (⋮) → \textbf{``Ver detalle''}. Se muestra la información completa: número de factura, cliente, fecha, ítems, totales y método de pago.

\subsection{Cancelar una venta}\label{cancelar-una-venta}

Una venta puede cancelarse mientras tenga estado \texttt{completed}:

\begin{enumerate}
\def\labelenumi{\arabic{enumi}.}
\tightlist
\item
  En la tabla de ventas, hacer clic en el menú de acciones (⋮) → \textbf{``Cancelar venta''}.
\item
  Ingresar el \textbf{motivo de cancelación} \emph{(requerido)}.
\item
  Confirmar.
\end{enumerate}

Al cancelar, el sistema:

\begin{itemize}
\tightlist
\item
  Cambia el estado de la venta a \texttt{cancelled}.
\item
  \textbf{Revierte el stock} de cada producto afectado (movimiento tipo \texttt{return}).
\item
  Registra la fecha y usuario de cancelación.
\end{itemize}

\subsection{Filtros y búsqueda}\label{filtros-y-buxfasqueda}

La tabla de ventas permite:

\begin{itemize}
\tightlist
\item
  Buscar por \textbf{número de factura}.
\item
  Filtrar por \textbf{estado} (\texttt{completed}, \texttt{cancelled}) y \textbf{método de pago}.
\item
  Ordenar por cualquier columna.
\end{itemize}

\section{Módulo de Clientes}\label{muxf3dulo-de-clientes}

El módulo permite gestionar el directorio de clientes de la organización y consultar su historial de compras.

\subsection{Registrar un cliente}\label{registrar-un-cliente}

\begin{enumerate}
\def\labelenumi{\arabic{enumi}.}
\tightlist
\item
  En el módulo Clientes, hacer clic en \textbf{``Agregar cliente''}.
\item
  Completar el formulario:

  \begin{itemize}
  \tightlist
  \item
    \textbf{Tipo de documento}: CC, NIT, Pasaporte, Otro \emph{(requerido)}.
  \item
    \textbf{Número de documento} \emph{(requerido, único por organización)}.
  \item
    \textbf{Nombre} y \textbf{apellido} \emph{(requeridos)}.
  \item
    \textbf{Correo electrónico}, \textbf{teléfono}, \textbf{dirección}, \textbf{ciudad} y \textbf{país} \emph{(opcionales)}.
  \item
    \textbf{Notas internas} \emph{(opcionales)}.
  \end{itemize}
\item
  Guardar.
\end{enumerate}

\subsection{Editar un cliente}\label{editar-un-cliente}

\begin{itemize}
\tightlist
\item
  Hacer clic en el menú de acciones (⋮) → \textbf{``Editar''}. Se pueden modificar todos los campos excepto el número de documento.
\end{itemize}

\subsection{Desactivar un cliente}\label{desactivar-un-cliente}

\begin{itemize}
\tightlist
\item
  Hacer clic en el menú de acciones (⋮) → \textbf{``Desactivar''}. El cliente se desactiva pero su historial de compras se conserva.
\end{itemize}

\subsection{Historial de compras del cliente}\label{historial-de-compras-del-cliente}

\begin{itemize}
\tightlist
\item
  Hacer clic en el menú de acciones (⋮) → \textbf{``Ver historial de compras''}.
\item
  Se muestran todas las ventas asociadas al cliente con fecha, número de factura, total y estado.
\item
  Se visualizan también los indicadores: \textbf{total acumulado de compras}, \textbf{número de compras} y \textbf{fecha de última compra}.
\end{itemize}

\section{Módulo de Administración}\label{muxf3dulo-de-administraciuxf3n}

\begin{quote}
Este módulo es exclusivo para usuarios con el rol \textbf{Administrador}.
\end{quote}

El módulo de administración permite gestionar las cuentas de usuario de la organización.

\subsection{Ver usuarios}\label{ver-usuarios}

La tabla muestra todos los usuarios de la organización con su nombre, correo, rol y estado (activo/inactivo).

\subsection{Crear un usuario}\label{crear-un-usuario}

\begin{enumerate}
\def\labelenumi{\arabic{enumi}.}
\tightlist
\item
  Hacer clic en \textbf{``Agregar usuario''}.
\item
  Completar el formulario:

  \begin{itemize}
  \tightlist
  \item
    \textbf{Nombre} y \textbf{apellido}.
  \item
    \textbf{Correo electrónico} --- será el identificador de ingreso.
  \item
    \textbf{Contraseña} temporal.
  \item
    \textbf{Rol} --- seleccionar entre los roles disponibles de la organización.
  \end{itemize}
\item
  Guardar.
\end{enumerate}

El nuevo usuario puede ingresar inmediatamente con las credenciales proporcionadas y cambiar su contraseña desde la configuración.

\subsection{Editar un usuario}\label{editar-un-usuario}

\begin{itemize}
\tightlist
\item
  Hacer clic en el menú de acciones (⋮) → \textbf{``Editar''}. Se pueden cambiar nombre, apellido y rol.
\end{itemize}

\subsection{Eliminar un usuario}\label{eliminar-un-usuario}

\begin{itemize}
\tightlist
\item
  Hacer clic en el menú de acciones (⋮) → \textbf{``Eliminar''}.
\item
  Confirmar la acción en el diálogo de confirmación.
\end{itemize}

\begin{quote}
\textbf{Advertencia:} Esta acción es permanente. El historial de ventas y movimientos del usuario permanece registrado.
\end{quote}

\section{Configuración}\label{configuraciuxf3n}

El módulo de configuración permite al usuario gestionar su cuenta personal y, si es administrador, los datos de la organización.

\subsection{Información personal}\label{informaciuxf3n-personal}

En la pestaña \textbf{``Perfil''}:

\begin{itemize}
\tightlist
\item
  Ver y editar \textbf{nombre} y \textbf{apellido}.
\item
  Ver el \textbf{correo electrónico} asociado (no editable desde aquí).
\end{itemize}

\subsection{Cambio de contraseña}\label{cambio-de-contraseuxf1a}

En la pestaña \textbf{``Seguridad''}:

\begin{enumerate}
\def\labelenumi{\arabic{enumi}.}
\tightlist
\item
  Ingresar la \textbf{contraseña actual}.
\item
  Ingresar la \textbf{nueva contraseña} y confirmarla.
\item
  Hacer clic en \textbf{``Actualizar contraseña''}.
\end{enumerate}

\subsection{Configuración de la organización}\label{configuraciuxf3n-de-la-organizaciuxf3n}

\begin{quote}
Disponible solo para usuarios con rol \textbf{Administrador}.
\end{quote}

En la pestaña \textbf{``Organización''}:

\begin{itemize}
\tightlist
\item
  Editar el \textbf{nombre} de la organización.
\item
  Editar el \textbf{slug} \emph{(identificador único)}.
\item
  Editar la \textbf{descripción} y la \textbf{URL del logo}.
\end{itemize}

\subsection{Apariencia}\label{apariencia}

En la pestaña \textbf{``Apariencia''}:

\begin{itemize}
\tightlist
\item
  Cambiar entre \textbf{modo claro} y \textbf{modo oscuro}.
\end{itemize}

\subsection{Eliminación de cuenta}\label{eliminaciuxf3n-de-cuenta}

En la pestaña \textbf{``Cuenta''}:

\begin{itemize}
\tightlist
\item
  El usuario puede solicitar la \textbf{eliminación de su propia cuenta}.
\item
  Se requiere confirmación escribiendo el correo electrónico.
\end{itemize}

\begin{quote}
\textbf{Nota:} La eliminación de cuenta no elimina el historial de transacciones asociadas.
\end{quote}

\section{Preguntas Frecuentes}\label{preguntas-frecuentes}

\textbf{¿Puedo usar OrbitEngine desde el celular?}\\
Sí. La interfaz está diseñada de manera responsiva y funciona en navegadores móviles modernos (Chrome, Safari). La experiencia óptima es en pantallas de escritorio o tabletas.

\textbf{¿Mis datos están separados de los de otras empresas?}\\
Sí. Cada organización tiene un espacio de trabajo completamente aislado. No es posible ver ni acceder a datos de otras organizaciones.

\textbf{¿Qué pasa si cancelo una venta? ¿El stock vuelve?}\\
Sí. Al cancelar una venta, el sistema revierte automáticamente el stock de todos los productos incluidos en ella.

\textbf{¿Puedo tener múltiples roles para usuarios?}\\
Actualmente cada usuario tiene un único rol. Los roles definen los permisos de acceso a los módulos y acciones disponibles.

\textbf{¿Cómo sé cuándo un producto está por agotarse?}\\
Los productos cuyo stock actual es igual o menor al \textbf{stock mínimo} configurado aparecen resaltados en la tabla de inventario y se contabilizan en el indicador ``Productos con stock bajo'' del dashboard.

\textbf{¿Se puede exportar la información a Excel o PDF?}\\
Los listados de inventario, ventas y clientes pueden exportarse a Excel (.xlsx) desde el módulo correspondiente, conservando los filtros aplicados. La exportación a PDF no está disponible en esta versión.

\textbf{¿Qué navegadores son compatibles?}\\
OrbitEngine es compatible con las versiones recientes de Chrome, Firefox, Edge y Safari. No se garantiza compatibilidad con Internet Explorer.

\emph{Documento generado como parte del proyecto de grado --- Universidad Sergio Arboleda, Semillero de Software como Innovación, Abril 2026.}
